\documentclass[a4paper,12pt,twoside]{article}
\usepackage[left=0.8in,right=0.8in,top=1in,bottom=1in]{geometry}

\usepackage{xeCJK}
\setCJKmainfont{Songti TC}
\setCJKfamilyfont{kai}{Songti TC}
\newcommand{\kai}{\CJKfamily{kai}} % 用 \kai{使用楷書}

\usepackage{lmodern}
\usepackage{titling}
\usepackage{titlesec}
\usepackage{graphicx}
\usepackage{float}
\usepackage{amsmath, amsthm, amssymb, amsfonts, mathrsfs}
\usepackage{upgreek} % 直立的希臘字母

\usepackage[shortlabels]{enumitem}
\usepackage{setspace}
\renewcommand{\baselinestretch}{1.5}
\usepackage[colorlinks]{hyperref}
\usepackage{fancyvrb}
\usepackage{fancyhdr}
\usepackage{multicol}
\usepackage{caption}
\usepackage{listings}
\usepackage{bm}
\usepackage{bbm} % for \mathbbm{1} (indicator function)
\usepackage{centernot}
\usepackage{mdframed}
\usepackage{tikz}
\usetikzlibrary{patterns, arrows.meta, positioning}
\usepackage[nameinlink, capitalise]{cleveref}
\usepackage{diagbox}
\usepackage{booktabs}
\usepackage{multirow,array}
\usepackage{tabularx}
\usepackage{mathtools}
\usepackage{indentfirst} % 首段 indent
\setlength{\parindent}{2em}
\usepackage{cancel} % 數學式子畫線槓掉
\newcommand\mcancel[2][red]{\setbox0=\hbox{$#2$}%
\rlap{\raisebox{.45\ht0}{\textcolor{#1}{\rule{\wd0}{1pt}}}}#2}
\usepackage{datetime2}

\AtBeginDocument{%
 \setlength{\abovedisplayskip}{6pt}
 \setlength{\belowdisplayskip}{6pt}
 \setlength{\abovedisplayshortskip}{4pt}
 \setlength{\belowdisplayshortskip}{4pt}
}

\definecolor{codegreen}{rgb}{0,0.6,0}
\definecolor{codegray}{rgb}{0.5,0.5,0.5}
\definecolor{codepurple}{rgb}{0.58,0,0.82}
\definecolor{backcolour}{rgb}{0.95,0.95,0.92}

\lstdefinestyle{mystyle}{
 backgroundcolor=\color{backcolour},
 commentstyle=\color{codegreen},
 keywordstyle=\color{magenta},
 numberstyle=\tiny\color{codegray},
 stringstyle=\color{codepurple},
 basicstyle=\ttfamily\footnotesize,
 breakatwhitespace=false,
 breaklines=true,
 captionpos=b,
 keepspaces=true,
 numbers=left,
 numbersep=5pt,
 showspaces=false,
 showstringspaces=false,
 showtabs=false,
 tabsize=2
}
\lstset{style=mystyle}

\hypersetup{citecolor=blue, linkcolor=blue, urlcolor=blue}

\DeclareMathOperator{\supp}{supp}
\DeclareMathOperator{\skewness}{Skewness}
\DeclareMathOperator{\DU}{D-U}
\DeclareMathOperator{\CU}{C-U}
\DeclareMathOperator{\Geo}{Geo}
\DeclareMathOperator{\Pois}{Pois}
\DeclareMathOperator{\Gam}{Gam}
\DeclareMathOperator{\LN}{LogN}
\DeclareMathOperator{\cov}{Cov}
\DeclareMathOperator*{\argmax}{arg\,max}
\DeclareMathOperator*{\argmin}{arg\,min}
\DeclareMathOperator{\re}{Re}
\DeclareMathOperator{\im}{Im}

\newcommand{\N}{\mathbb{N}}
\newcommand{\Z}{\mathbb{Z}}
\newcommand{\Q}{\mathbb{Q}}
\newcommand{\R}{\mathbb{R}}
\newcommand{\C}{\mathbb{C}}
\newcommand{\G}{\mathcal{G}}
\newcommand{\mP}{\mathbb{P}}
\newcommand{\F}{\mathcal{F}}
\newcommand{\B}{\mathcal{B}}
\newcommand{\1}{\mathbbm{1}}
\newcommand{\Poi}{\text{Poi}}
\newcommand{\Var}{\text{Var}}
\newcommand{\as}{\text{a.s.}}
\newcommand{\vect}[1]{\boldsymbol{\mathbf{#1}}}
\newcommand{\paren}[1]{\left({#1}\right)}
\newcommand{\brk}[1]{\left[{#1}\right]}
\newcommand{\curl}[1]{\left\{{#1}\right\}}
\renewcommand{\lambda}{\lambdaup}

\theoremstyle{definition} % 定理環境 style 改用正體
\newtheorem{theorem}{定理}
\newtheorem{definition}{定義}
\newtheorem*{definition*}{定義}
\newtheorem{example}{例子}
\newtheorem{exercise}{練習}
\newtheorem{corollary}{推論}
\newtheorem{property}{性質}
\newtheorem{prop}{命題}
\newtheorem{lemma}{引理}
\newtheorem{problem}{問題}
\newtheorem*{remark}{註}

\renewcommand{\figurename}{圖}
\renewcommand{\tablename}{表}

\DeclareMathOperator{\var}{Var}
\DeclareMathOperator{\corr}{corr}
\DeclareMathOperator{\rank}{rank}
\DeclareMathOperator{\MSE}{MSE}
\DeclareMathOperator{\SSE}{SSE}
\DeclareMathOperator{\SST}{SST}
\DeclareMathOperator{\SSR}{SSR}
\DeclareMathOperator{\sgn}{sign}
\DeclareMathOperator{\alt}{Alt}
\DeclareMathOperator{\adj}{adj}
\DeclareMathOperator{\curlOp}{curl}
\DeclareMathOperator{\spn}{span}
\DeclareMathOperator{\sse}{SSE}
\DeclareMathOperator{\tr}{tr}

\newcommand{\ep}{\varepsilon}
\newcommand{\p}{\partial}
\newcommand{\abs}[1]{\left|#1\right|}
\newcommand\norm[1]{\left\lVert#1\right\rVert}
\newcommand{\innerprod}[1]{\langle#1\rangle}
\newcommand{\ds}{\displaystyle}
\newcommand{\blank}{\rule[-1mm]{1.5cm}{0.1mm}}
\newcommand{\dint}{\displaystyle\int}
\renewcommand{\bar}{\overline}

\newtheoremstyle{remark}{}{}{}{}{\bfseries\sffamily}{.}{.5em}{#1}
\theoremstyle{remark}

\newtheorem*{remark*}{\textit{Remark}}
\newtheorem*{solution*}{解}

\AtEndEnvironment{solution*}{\hfill$\square$}
\newtheorem*{hint*}{Hint}


\usepackage{diagbox}

\fancypagestyle{firstpage}{
  \fancyhf{}
  \fancyhead[R]{\small《高中數學實驗教材》第 5 冊（自然組）}
  \fancyhead[L]{\small 附錄\quad 階乘、二項係數表}
  \fancyfoot[C]{\thepage}
}

\pagestyle{fancy}
\setlength{\headheight}{14pt}
\renewcommand{\headrulewidth}{0.5pt}
\fancyhead{}
\fancyhead[LE]{\small 《高中數學實驗教材》第五冊（自然組）}
\fancyhead[RO]{\small 附錄\quad 階乘、二項係數表}
\fancyfoot[C]{\thepage}

\title{\LARGE\textbf{《高中數學實驗教材》第五冊（自然組）}\\[0.5em] \Large\textbf{附錄\quad 階乘、二項係數表}}
\author{\textbf{高中數學實驗教材編輯小組}\quad 編著}
\date{民國六十四年（1975）再版}

\begin{document}

\maketitle
\thispagestyle{firstpage}

% ==============================================================================
% 《高中數學實驗教材》第五冊（自然組）
% 附錄：階乘、二項係數表
% ==============================================================================

\section*{附錄\quad 階乘、二項係數表}
\addcontentsline{toc}{section}{附錄\texorpdfstring{\quad}{ }階乘、二項係數表}
\label{sec:b5-appendix}
\markboth{附錄\quad 階乘、二項係數表}{附錄\quad 階乘、二項係數表}

\vspace{-0.5em}

\begin{center}
{\large\textbf{階\qquad 乘}}
\end{center}

\vspace{-0.3em}

\noindent\begin{minipage}[t]{0.68\textwidth}
\centering\fontsize{7.5pt}{9pt}\selectfont\renewcommand{\arraystretch}{1.0}\setlength{\tabcolsep}{3.5pt}
\begin{tabular}{|r|r|l|}
\hline
\multicolumn{1}{|c|}{$n$} & \multicolumn{1}{c|}{$n!$} & \multicolumn{1}{c|}{$1/n!$} \\ \hline
1 & 1 & 1 \\
2 & 2 & 0.5 \\
3 & 6 & 0.166666666666666666 \\
4 & 24 & 0.041666666666666666 \\
5 & 120 & 0.008333333333333333 \\
6 & 720 & 0.001388888888888888 \\
7 & 5040 & 0.000198412698412698 \\
8 & 40320 & 0.000024801587301587 \\
9 & 362880 & 0.000002755731922398 \\
10 & 3628800 & 0.000000275573192239 \\
11 & 39916800 & 0.000000025052108385 \\
12 & 479001600 & 0.000000002087675698 \\
13 & 6227020800 & 0.000000000160590438 \\
14 & 87178291200 & 0.000000000011470745 \\
15 & 1307674368000 & 0.000000000000764716 \\
16 & 20922789888000 & 0.000000000000047794 \\
17 & 355687428096000 & 0.000000000000002811 \\
18 & 6402373705728000 & 0.000000000000000156 \\
19 & 121645100408832000 & $0.82206 \times 10^{-17}$ \\
20 & 2432902008176640000 & $0.41103 \times 10^{-18}$ \\
\hline
\end{tabular}
\end{minipage}\hfill
\begin{minipage}[t]{0.30\textwidth}
\centering\fontsize{7.5pt}{9pt}\selectfont\renewcommand{\arraystretch}{1.0}\setlength{\tabcolsep}{5pt}
\begin{tabular}{|r|c|r|}
\hline
\multicolumn{3}{|c|}{$n! = a \times 10^c$} \\ \hline
\multicolumn{1}{|c|}{$n$} & $a$ & \multicolumn{1}{c|}{$c$} \\ \hline
25 & 1.5511 & 25 \\
30 & 2.6525 & 32 \\
35 & 1.0333 & 40 \\
40 & 8.1592 & 47 \\
45 & 1.1962 & 56 \\
50 & 3.0414 & 64 \\
60 & 8.3210 & 81 \\
70 & 1.1979 & 100 \\
80 & 7.1569 & 118 \\
90 & 1.4857 & 138 \\
100 & 9.3326 & 157 \\
120 & 6.6895 & 198 \\
140 & 1.3462 & 241 \\
160 & 4.7147 & 284 \\
180 & 2.0090 & 329 \\
200 & 7.8866 & 374 \\
220 & 2.2839 & 421 \\
240 & 4.0679 & 468 \\ \hline
\end{tabular}
\end{minipage}

\vspace{0.8em}

\begin{center}
{\large\textbf{二\quad 項\quad 係\quad 數}}\\[0.3em]
$\displaystyle \binom{N}{n} = \frac{N!}{n!(N-n)!}$
\end{center}

\vspace{-0.3em}

\begin{center}
\fontsize{7pt}{8.2pt}\selectfont\setlength{\tabcolsep}{2pt}\renewcommand{\arraystretch}{0.96}
\begin{tabular}{|c|cccccccccccccccccccc|}
\hline
\diagbox[width=2.5em,height=1.7em]{$n$}{$N$} & 1 & 2 & 3 & 4 & 5 & 6 & 7 & 8 & 9 & 10 & 11 & 12 & 13 & 14 & 15 & 16 & 17 & 18 & 19 & 20 \\ \hline
1 & 1 & 2 & 3 & 4 & 5 & 6 & 7 & 8 & 9 & 10 & 11 & 12 & 13 & 14 & 15 & 16 & 17 & 18 & 19 & 20 \\
2 &  & 1 & 3 & 6 & 10 & 15 & 21 & 28 & 36 & 45 & 55 & 66 & 78 & 91 & 105 & 120 & 136 & 153 & 171 & 190 \\
3 &  &  & 1 & 4 & 10 & 20 & 35 & 56 & 84 & 120 & 165 & 220 & 286 & 364 & 455 & 560 & 680 & 816 & 969 & 1140 \\
4 &  &  &  & 1 & 5 & 15 & 35 & 70 & 126 & 210 & 330 & 495 & 715 & 1001 & 1365 & 1820 & 2380 & 3060 & 3876 & 4845 \\
5 &  &  &  &  & 1 & 6 & 21 & 56 & 126 & 252 & 462 & 792 & 1287 & 2002 & 3003 & 4368 & 6188 & 8568 & 11628 & 15504 \\
6 &  &  &  &  &  & 1 & 7 & 28 & 84 & 210 & 462 & 924 & 1716 & 3003 & 5005 & 8008 & 12376 & 18564 & 27132 & 38760 \\
7 &  &  &  &  &  &  & 1 & 8 & 36 & 120 & 330 & 792 & 1716 & 3432 & 6435 & 11440 & 19448 & 31824 & 50388 & 77520 \\
8 &  &  &  &  &  &  &  & 1 & 9 & 45 & 165 & 495 & 1287 & 3003 & 6435 & 12870 & 24310 & 43758 & 75582 & 125970 \\
9 &  &  &  &  &  &  &  &  & 1 & 10 & 55 & 220 & 715 & 2002 & 5005 & 11440 & 24310 & 48620 & 92378 & 167960 \\
10 &  &  &  &  &  &  &  &  &  & 1 & 11 & 66 & 286 & 1001 & 3003 & 8008 & 19448 & 43758 & 92378 & 184756 \\
11 &  &  &  &  &  &  &  &  &  &  & 1 & 12 & 78 & 364 & 1365 & 4368 & 12376 & 31824 & 75582 & 167960 \\
12 &  &  &  &  &  &  &  &  &  &  &  & 1 & 13 & 91 & 455 & 1820 & 6188 & 18564 & 50388 & 125970 \\
13 &  &  &  &  &  &  &  &  &  &  &  &  & 1 & 14 & 105 & 560 & 2380 & 8568 & 27132 & 77520 \\
14 &  &  &  &  &  &  &  &  &  &  &  &  &  & 1 & 15 & 120 & 680 & 3060 & 11628 & 38760 \\
15 &  &  &  &  &  &  &  &  &  &  &  &  &  &  & 1 & 16 & 136 & 816 & 3876 & 15504 \\
16 &  &  &  &  &  &  &  &  &  &  &  &  &  &  &  & 1 & 17 & 153 & 969 & 4845 \\
17 &  &  &  &  &  &  &  &  &  &  &  &  &  &  &  &  & 1 & 18 & 171 & 1140 \\
18 &  &  &  &  &  &  &  &  &  &  &  &  &  &  &  &  &  & 1 & 19 & 190 \\
19 &  &  &  &  &  &  &  &  &  &  &  &  &  &  &  &  &  &  & 1 & 20 \\
20 &  &  &  &  &  &  &  &  &  &  &  &  &  &  &  &  &  &  &  & 1 \\
\hline
\end{tabular}
\end{center}


\end{document}
